\documentclass[conference]{IEEEtran}

\usepackage[utf8]{inputenc}
\usepackage[T1]{fontenc}
\usepackage[spanish,es-nodecimaldot]{babel}
\usepackage{amsmath,amssymb,mathtools}
\usepackage{siunitx}
\usepackage{graphicx}
\usepackage{booktabs}
\usepackage{array}
\usepackage{xcolor}
\usepackage{microtype}
\usepackage{cite}
\usepackage{flushend}
\usepackage{listings}
\usepackage{tikz}
\usetikzlibrary{babel,decorations.pathmorphing}
\usepackage[hidelinks]{hyperref}

\sisetup{per-mode=symbol}
\renewcommand{\IEEEkeywordsname}{Palabras clave}

\definecolor{draftgray}{RGB}{244,246,249}
\definecolor{draftblue}{RGB}{31,77,120}

\newcommand{\draftnote}[1]{%
  \par\smallskip
  \noindent\colorbox{draftgray}{%
    \parbox{\dimexpr\columnwidth-2\fboxsep\relax}{%
      \footnotesize\color{draftblue}\textbf{Pendiente de completar:} #1}}
  \par\smallskip
}

\lstdefinelanguage{Julia}{
  morekeywords={abstract,type,immutable,bitstype,typealias,function,end,quote,
  if,else,elseif,while,for,begin,let,do,try,catch,finally,return,break,
  continue,macro,ccall,local,global,const,struct,mutable,where,module,using,
  import,export,primitive},
  sensitive=true,
  morecomment=[l]{\#},
  morestring=[b]",
  morestring=[m]'
}

\lstset{
  language=Julia,
  basicstyle=\ttfamily\scriptsize,
  keywordstyle=\color{blue!60!black}\bfseries,
  commentstyle=\color{green!40!black},
  stringstyle=\color{red!55!black},
  numbers=left,
  numberstyle=\tiny\color{gray},
  stepnumber=1,
  numbersep=6pt,
  frame=single,
  breaklines=true,
  columns=fullflexible,
  showstringspaces=false,
  tabsize=2
}

\title{Modelado y simulación de un sistema\\masa--resorte--amortiguador aplicado a una rueda}

\author{
  \IEEEauthorblockN{Maria José Álvarez Hernández \qquad Cristian Alejandro Guarin Marin}
  \IEEEauthorblockA{
    Matemáticas Especiales\\
    Universidad de Antioquia\\
    Medellín, Colombia}
}

\begin{document}

\maketitle

\begin{abstract}
Este trabajo formula un modelo lineal de un grado de libertad para representar el movimiento vertical de la masa equivalente asociada a una rueda conectada a un chasis fijo mediante un resorte y un amortiguador viscoso. La ecuación de movimiento se obtiene mediante el formalismo de Lagrange con disipación de Rayleigh y se expresa posteriormente como una función de transferencia entre una fuerza perturbadora del terreno y el desplazamiento respecto al equilibrio estático. A partir de la forma estándar de segundo orden se identifican la frecuencia natural y la razón de amortiguamiento. El modelo servirá como base para evaluar en Julia sus respuestas natural y forzada.
\end{abstract}

\begin{IEEEkeywords}
sistema masa--resorte--amortiguador, ecuaciones de Lagrange, función de transferencia, respuesta dinámica, Julia.
\end{IEEEkeywords}

\draftnote{Reemplazar los nombres de las integrantes y actualizar el resumen cuando existan parámetros, resultados numéricos y conclusiones definitivas.}

\section{Introducción}

Los sistemas masa--resorte--amortiguador constituyen una aproximación fundamental para estudiar vibraciones mecánicas y respuestas transitorias. En este trabajo se emplea esta estructura para construir una analogía simplificada del movimiento vertical de una rueda: la masa equivalente se conecta a un chasis considerado fijo mediante un elemento elástico y un amortiguador viscoso, mientras que las irregularidades del terreno se representan mediante una fuerza externa variable en el tiempo.

El objetivo es obtener el modelo matemático mediante las ecuaciones de Lagrange, expresarlo en el dominio de Laplace y evaluar su comportamiento dinámico mediante un programa en Julia. Se considerarán dos experimentos: una respuesta natural desde una condición inicial no nula y una respuesta forzada desde el equilibrio y el reposo. Esta separación permitirá observar tanto las propiedades intrínsecas del sistema como su reacción ante una perturbación externa.

El modelo no pretende reproducir la dinámica completa de un vehículo. Su alcance se limita a un grado de libertad y omite, entre otros efectos, el movimiento del chasis, la rotación de la rueda, la elasticidad independiente del neumático y las no linealidades propias de una suspensión real.

\section{Descripción y convenciones del modelo}

La Fig.~\ref{fig:sistema} muestra el esquema adoptado. La coordenada generalizada $y(t)$ describe el desplazamiento vertical de la masa equivalente respecto a su posición de equilibrio estático y se considera positiva hacia arriba. Los extremos superiores del resorte y del amortiguador se conectan a un chasis fijo; sus extremos inferiores actúan sobre el eje asociado a la masa equivalente.

\begin{figure}[t]
  \centering
  \begin{tikzpicture}[
  scale=0.78,
  every node/.style={font=\footnotesize},
  spring/.style={decorate,decoration={zigzag,segment length=3.2mm,amplitude=1.5mm}}
]
  % Chasís fijo
  \draw[very thick] (-1.65,1.90) -- (1.65,1.90);
  \foreach \x in {-1.5,-1.2,...,1.5}
    \draw[thin] (\x,1.90) -- ++(0.18,0.18);
  \node[anchor=south] at (0,2.15) {Chasís fijo};

  % Rueda y eje
  \draw[very thick] (0,0) circle (1.35);
  \draw[thick] (0,0) circle (0.95);
  \draw[fill=white,thick] (0,0) circle (0.16);
  \node[anchor=north] at (0,-0.18) {$m$};

  % Línea de equilibrio
  \draw[densely dashed,gray] (-1.52,0) -- (1.52,0);
  \node[anchor=east,gray] at (-1.52,0) {$y=0$};

  % Amortiguador
  \draw[thick] (-0.50,1.90) -- (-0.50,1.52);
  \draw[thick] (-0.72,1.26) rectangle (-0.28,1.52);
  \draw[thick] (-0.66,1.38) -- (-0.34,1.38);
  \draw[thick] (-0.50,1.38) -- (-0.50,0.28);
  \draw[thick] (-0.50,0.28) -- (0,0);
  \node[anchor=east] at (-0.72,1.38) {$\beta$};

  % Resorte
  \draw[thick] (0.50,1.90) -- (0.50,1.67);
  \draw[thick,spring] (0.50,1.67) -- (0.50,0.35);
  \draw[thick] (0.50,0.35) -- (0,0);
  \node[anchor=west] at (0.68,1.05) {$k$};

  % Suelo y fuerza
  \draw[very thick] (-1.55,-1.42) -- (1.55,-1.42);
  \draw[->,very thick] (0,-2.02) -- (0,-1.45)
    node[midway,right] {$f(t)$};

  % Coordenada positiva
  \draw[->,thick] (1.72,-0.28) -- (1.72,0.78)
    node[above] {$+y$};
\end{tikzpicture}


  \caption{Modelo simplificado de la rueda con resorte, amortiguador y fuerza perturbadora del terreno.}
  \label{fig:sistema}
\end{figure}

La entrada $f(t)$ no representa la fuerza normal total, sino su variación respecto al valor estático. Si $N(t)$ es la fuerza normal total y $N_0$ su valor de equilibrio, entonces
\begin{equation}
  N(t)=N_0+f(t).
\end{equation}
En el equilibrio, la fuerza normal estática compensa el peso de la masa equivalente y la carga estática transmitida por la suspensión. Por ello, al medir $y(t)$ desde dicha posición, las contribuciones constantes se cancelan y no aparecen en la ecuación dinámica incremental.

\begin{table}[t]
  \caption{Variables y parámetros del modelo}
  \label{tab:simbolos}
  \centering
  \begin{tabular}{@{}cll@{}}
    \toprule
    Símbolo & Descripción & Unidad \\
    \midrule
    $m$ & Masa equivalente de la rueda & \si{\kilogram} \\
    $k$ & Rigidez equivalente del resorte & \si{\newton\per\meter} \\
    $\beta$ & Coeficiente de amortiguamiento & \si{\newton\second\per\meter} \\
    $y(t)$ & Desplazamiento desde el equilibrio & \si{\meter} \\
    $f(t)$ & Fuerza perturbadora del terreno & \si{\newton} \\
    \bottomrule
  \end{tabular}
\end{table}

El desarrollo supone movimiento exclusivamente vertical, parámetros constantes, comportamiento lineal del resorte y amortiguamiento viscoso lineal. También se desprecia la masa propia del resorte y del amortiguador.

\section{Derivación mediante las ecuaciones de Lagrange}

El sistema posee un único grado de libertad, de modo que la coordenada generalizada se define como
\begin{equation}
  q(t)=y(t).
\end{equation}
La masa equivalente se traslada verticalmente con velocidad $\dot y(t)$, por lo que su energía cinética es
\begin{equation}
  T=\frac{1}{2}m\dot y^{,2}.
  \label{eq:cinetica}
\end{equation}
La deformación dinámica del resorte se mide desde el equilibrio y su energía potencial elástica corresponde a
\begin{equation}
  V=\frac{1}{2}ky^2.
  \label{eq:potencial}
\end{equation}
No se incluye un término gravitacional independiente porque la coordenada se ha definido respecto al equilibrio estático. Incluir únicamente $mgy$, sin incorporar las demás fuerzas constantes que lo compensan, haría que $y=0$ dejara de representar un equilibrio.

El lagrangiano se obtiene como la diferencia entre las energías cinética y potencial:
\begin{equation}
  L=T-V=\frac{1}{2}m\dot y^{,2}-\frac{1}{2}ky^2.
  \label{eq:lagrangiano}
\end{equation}

El amortiguador no almacena energía potencial; disipa energía mediante una fuerza opuesta a la velocidad. Para un amortiguador viscoso lineal se emplea la función de disipación de Rayleigh
\begin{equation}
  \mathcal{D}=\frac{1}{2}\beta\dot y^{,2}.
  \label{eq:rayleigh}
\end{equation}
Como la fuerza perturbadora actúa en el mismo sentido positivo de la coordenada, la fuerza generalizada es $Q_y=f(t)$. Esto también se obtiene del trabajo virtual $\delta W=f(t)\,\delta y$.

La ecuación de Lagrange con disipación es
\begin{equation}
  \frac{\mathrm{d}}{\mathrm{d}t}
  \left(\frac{\partial L}{\partial \dot y}\right)
  -\frac{\partial L}{\partial y}
  +\frac{\partial\mathcal{D}}{\partial\dot y}
  =Q_y.
  \label{eq:lagrange}
\end{equation}
Las derivadas necesarias son
\begin{align}
  \frac{\partial L}{\partial\dot y} &= m\dot y,
  &\frac{\mathrm{d}}{\mathrm{d}t}
  \left(\frac{\partial L}{\partial\dot y}\right) &= m\ddot y,\\
  \frac{\partial L}{\partial y} &= -ky,
  &\frac{\partial\mathcal D}{\partial\dot y} &= \beta\dot y.
\end{align}
Al sustituirlas en \eqref{eq:lagrange} se obtiene
\begin{equation}
  \boxed{m\ddot y(t)+\beta\dot y(t)+ky(t)=f(t)}.
  \label{eq:movimiento}
\end{equation}
La fuerza restauradora del resorte puede reconocerse al despejar la aceleración: aparece como $-ky$, mientras que la fuerza del amortiguador aparece como $-\beta\dot y$. Ambas se oponen, respectivamente, al desplazamiento y a la velocidad.


\section{Función de transferencia}

Se define $f(t)$ como la entrada y $y(t)$ como la salida. Para obtener la función de transferencia se aplica la transformada de Laplace a \eqref{eq:movimiento} bajo condiciones iniciales nulas,
\begin{equation}
  y(0)=0, \qquad \dot y(0)=0.
\end{equation}
En consecuencia,
\begin{equation}
  ms^2Y(s)+\beta sY(s)+kY(s)=F(s).
\end{equation}
Factorizando $Y(s)$ resulta
\begin{equation}
  Y(s)\left(ms^2+\beta s+k\right)=F(s),
\end{equation}
y la función de transferencia entre la fuerza perturbadora y el desplazamiento es
\begin{equation}
  \boxed{G(s)=\frac{Y(s)}{F(s)}
  =\frac{1}{ms^2+\beta s+k}}.
  \label{eq:transferencia}
\end{equation}

Para reconocer la estructura estándar de segundo orden se divide la ecuación de movimiento entre $m$:
\begin{equation}
  \ddot y+\frac{\beta}{m}\dot y+\frac{k}{m}y
  =\frac{1}{m}f(t).
  \label{eq:normalizada}
\end{equation}
Al comparar su lado izquierdo con
\begin{equation}
  \ddot y+2\zeta\omega_n\dot y+\omega_n^2y,
\end{equation}
se identifican
\begin{equation}
  \boxed{\omega_n=\sqrt{\frac{k}{m}}},
  \qquad
  \boxed{\zeta=\frac{\beta}{2\sqrt{mk}}}.
  \label{eq:parametros_dinamicos}
\end{equation}
Por tanto, \eqref{eq:transferencia} también puede escribirse como
\begin{equation}
  G(s)=
  \frac{\dfrac{1}{k}\omega_n^2}
  {s^2+2\zeta\omega_ns+\omega_n^2}.
  \label{eq:transferencia_estandar}
\end{equation}
Esta forma muestra que la ganancia estática es $G(0)=1/k$ y que el denominador determina la frecuencia y el amortiguamiento de la respuesta transitoria.


\section{Caracterización dinámica y parámetros}

El comportamiento libre está gobernado por la ecuación característica
\begin{equation}
  ms^2+\beta s+k=0,
\end{equation}
cuyas raíces son
\begin{equation}
  s_{1,2}=\frac{-\beta\pm\sqrt{\beta^2-4mk}}{2m}.
  \label{eq:polos}
\end{equation}
Según el valor de $\beta^2-4mk$, el sistema será subamortiguado, críticamente amortiguado o sobreamortiguado. En términos de $\zeta$, estos regímenes corresponden a $0<\zeta<1$, $\zeta=1$ y $\zeta>1$, respectivamente. Si $m>0$, $\beta>0$ y $k>0$, los polos presentan parte real negativa y el equilibrio es asintóticamente estable.

Para el caso subamortiguado se define la frecuencia natural amortiguada
\begin{equation}
  \omega_d=\omega_n\sqrt{1-\zeta^2}.
\end{equation}
La respuesta natural tendrá entonces la forma general
\begin{equation}
  y_n(t)=e^{-\zeta\omega_nt}
  \left[A\cos(\omega_dt)+B\sin(\omega_dt)\right],
  \label{eq:respuesta_natural}
\end{equation}
donde $A$ y $B$ se determinan a partir de las condiciones iniciales.

\draftnote{Seleccionar y justificar los valores de $m$, $k$ y $\beta$; calcular $\omega_n$, $\zeta$, $\omega_d$ y los polos; completar una tabla con símbolo, valor, unidad, fuente y criterio físico.}


\section{Metodología de simulación y evaluación}

La evaluación se dividirá en dos casos. La respuesta natural caracterizará el comportamiento intrínseco del sistema sin excitación externa,
\begin{equation}
  f(t)=0, \qquad y(0)=y_0\neq0, \qquad \dot y(0)=0.
\end{equation}
Aunque esta condición representa una rueda desplazada y liberada, su inclusión permite observar directamente el decaimiento asociado al amortiguamiento.

La respuesta forzada constituirá el escenario principal de la analogía vehicular. El sistema comenzará en equilibrio y reposo,
\begin{equation}
  y(0)=0, \qquad \dot y(0)=0,
\end{equation}
y recibirá una fuerza perturbadora conocida que represente una irregularidad del terreno. Una entrada candidata es un pulso suave de media onda senoidal,
\begin{equation}
  f(t)=
  \begin{cases}
    F_0\sin\left(\dfrac{\pi t}{T_p}\right), & 0\leq t\leq T_p,\\[5pt]
    0, & t>T_p,
  \end{cases}
  \label{eq:pulso}
\end{equation}
donde $F_0$ es la fuerza máxima y $T_p$ su duración. La selección definitiva de la entrada deberá conservarse de manera idéntica en la formulación matemática, el programa y las figuras.

\draftnote{Documentar los paquetes de Julia, el método de simulación, el intervalo temporal, las tolerancias y el procedimiento reproducible para generar y guardar las figuras.}


\newpage
\section{Resultados y comparación}

La respuesta natural deberá contrastarse con la frecuencia amortiguada y con el factor de decaimiento $e^{-\zeta\omega_nt}$. Para la respuesta forzada se verificará que el desplazamiento parta de cero, responda mientras actúa la perturbación y regrese al equilibrio una vez desaparezca la fuerza, siempre que los polos se encuentren en el semiplano izquierdo.

La comparación considerará, como mínimo, la amplitud máxima, el número de oscilaciones visibles, el tiempo de decaimiento y el comportamiento final. También deberá distinguirse la respuesta provocada por energía inicial de aquella producida por la fuerza externa.

\draftnote{Insertar aquí las figuras exportadas por Julia, sus parámetros, las verificaciones numéricas y la interpretación física. Cada figura debe incluir título, ejes, unidades y leyenda.}


\section{Conclusión}

\draftnote{Redactar una única conclusión de menos de 100 palabras después de comparar las respuestas natural y forzada. Evitar introducir resultados nuevos en esta sección.}



\appendices
\section{Código completo en Julia}

\draftnote{Incluir exactamente la versión del código que genere las figuras definitivas. El anexo deberá leer directamente el archivo \texttt{code/assignment1.jl} para mantenerse sincronizado con el programa ejecutado.}

% Descomentar cuando exista el archivo definitivo:
% \lstinputlisting[caption={Programa completo en Julia},label={lst:julia}]{code/assignment1.jl}


% Activar cuando se incorporen fuentes bibliográficas:
% \bibliographystyle{IEEEtran}
% \bibliography{references}

\end{document}
